\chapter{Discussion}
\label{ch:6}

\section{Answers to the Research Questions}
\label{sec:6.1}
\begin{description}
  \item[RQ1. Does multi-teacher distillation with per-instance weighting improve the compact student?]
    Not in sample. Across five students, three committees and every weighting temperature from 0.05 to 5,
    no distilled student beats the same student fine-tuned without teachers by more than the test set can
    resolve, and per-instance weighting never differs from uniform averaging (Sections~\ref{sec:5.2}
    and~\ref{sec:5.3}). A teacher trained on labelled implication changes nothing either
    (Section~\ref{sec:5.5}).
  \item[RQ2. Why not?] The student is distilled on the split the teachers were fine-tuned on, where they
    have memorised the labels. There their soft labels are the gold labels, the reliability signal that
    weights them is saturated, and task adaptation has made them agree on more than nine tweets in ten
    (Sections~\ref{sec:5.4} and~\ref{sec:5.6}).
  \item[RQ3. Does distillation on unlabelled text help, and on what does the gain depend?] Yes. On
    BERT-mini the gain grows monotonically with the amount of unlabelled text, from +0.0010 at 5,000
    tweets to +0.0123 [+0.0035, +0.0211] at 168,000. Generic tweets from the same platform work nearly as
    well as abuse-related ones, fine-tuning for the same number of updates gains nothing, and the gain
    holds on the BiLSTM (+0.0163) and on BERT-small (+0.0057, with an interval that includes zero). One
    teacher is enough: the committee, its weighting and hard pseudo-labels add nothing
    (Sections~\ref{sec:5.7} and~\ref{sec:5.8}).
  \item[RQ4. Does any of it improve the discrimination of implied abuse from harmless sarcasm?] No. Every
    pre-trained variant of the student, in sample or out, ranks ironic abuse above benign sarcasm with an
    AUC between 0.73 and 0.78. The interventions change how readily the student fires, not what it can
    tell apart (Section~\ref{sec:5.9}).
\end{description}

\section{What the Findings Mean}
\label{sec:6.2}

\textbf{What the audit says about error-weighted multi-teacher distillation.} The family's weighting
rule scores each teacher's reliability against the gold label on the training set. For teachers
fine-tuned on that set the score is saturated, and the weights it produces are uniform whatever the
temperature. Correcting the estimate, by scoring reliability out of sample as Section~\ref{sec:3.5}
does or by fitting a gate on held-out data, is necessary for the rule to mean anything, and on a
committee assembled in the usual way it is not sufficient: the teachers' outputs do not carry the
information needed to pick the right one where they disagree, so the reachable gain over uniform
averaging is a few tenths of a point, in sample and out. The five-point oracle gap is real, and it
is made of instances on which the teachers, and very likely the annotators, are uncertain.

\textbf{What it says about the committee.} Assembling specialists for complementary expertise and
then fine-tuning them on one corpus produces teachers that agree at kappa 0.9 and share their
errors. The adaptation cannot be skipped, because it is what makes the specialists' outputs
comparable; it can only be paid for. A committee that keeps its diversity needs at least one member
whose knowledge was not obtained by fitting the target training set. Out of sample, where the
committee does hold a better label than any of its members, the student still receives no more from
three teachers than from one: on this task a single large teacher's soft labels on unseen text are
the whole of what distillation transfers, and the cost of a committee is three times the caching for
nothing.

\textbf{What it says about where distillation's gain is.} In sample the teachers' soft labels are
the gold labels and distillation reduces to fine-tuning; out of sample they carry information the
gold labels do not, and a compact student receives it in proportion to how much such text it sees.
The recipes that report large gains for small students distil on transfer sets larger than the
labelled data \cite{tang2019distilling, jiao2020tinybert, turc2019wellread}, and Beyer et al.'s
account of distillation as function matching says why: a function is learned where it is sampled,
and sampling it only where it equals the label teaches the label \cite{beyer2022knowledge}. What our
measurements add is the domain-specific shape. The curve is monotone over a 34-fold range of size
and has not saturated at 168,000 rows; generic tweets carry most of the value, so the text need not
be abuse-related, only in the register the student will meet; the gain is not optimisation length,
which on its own overfits from the thirteenth epoch; and the heterogeneous student, which cannot
share the teacher's representation, gains most, which places the effect in the soft labels rather
than in the hidden-state term. The cost is one forward pass of one teacher over the unlabelled text
and a longer student run; the deployed model is unchanged.

\textbf{What it says about implication.} The gain lands on the two classes the labelled split draws
worst, not\_cyberbullying and other\_cyberbullying, and nowhere else. Recall at a fixed threshold
and the false-positive rate on harmless sarcasm move together across 184 models; a method that
reports only the first can claim progress on implication by lowering its threshold, and the
out-of-sample students do the opposite, firing less on sarcasm of both kinds while discriminating no
better. The threshold-free measure shows the ability to be fixed by pre-training and unmoved by
every objective, committee, teacher and transfer set we tried. If unlabelled text is to teach
implication, its composition is the variable, not its size: nothing in 205,593 tweets of offence and
sentiment carries the distinction the probes measure, and a transfer set that did would have to be
built from text in which implication is present and marked, which is the corpus of
Section~\ref{sec:4.2} and not the platform's ordinary stream.

\textbf{What it says about reporting.} Every out-of-sample result in this paper was predicted in
writing before the run that tested it, with thresholds and a decision rule, and scored by a script
against the generated tables; three of the eleven predictions failed as written and are reported as
such. A test set of 4,326 posts cannot resolve differences under 0.007 macro-F1, which is larger
than every in-sample distillation effect in the grid and than most differences reported in this
literature. The multi-teacher results the field has published for this task were obtained without
the fine-tuning control, the teachers' own scores, seeds or intervals; under those conditions our
in-sample grid would also have read as a success.

\section{How the Study Changed}
\label{sec:6.3}
The design of this thesis changed four times, and each change was forced by a measurement rather than
chosen in advance. Every change is recorded, with the evidence that forced it, in the decision log of
the project's repository.
\begin{enumerate}
  \item \textbf{From the Phase-2 method to a controlled grid.} The Pre-Thesis II report distilled two
    teachers into DistilBERT on the Wikipedia corpus and reported the student above its teachers from one
    run, without a student trained without teachers. On the cleaned tweet corpus we added that control,
    three seeds and paired bootstrap intervals, and distillation then raised every student by less than
    the test set can resolve.
  \item \textbf{From dynamic weighting to its mechanism.} The per-instance weights sat close to uniform at
    the default temperature and barely sharpened at a twentieth of it, and the score did not move.
    Measuring the teachers explained why: they had memorised the training split and agreed with one
    another on more than nine tweets in ten.
  \item \textbf{From a specialist teacher to the implication audit.} A teacher trained on labelled
    implication was added to supply the missing expertise. After task adaptation it behaved almost exactly
    like HateBERT, and no student became better at telling implied abuse from harmless sarcasm; the
    threshold-free measure showed that ability fixed by pre-training.
  \item \textbf{From the audit to out-of-sample distillation.} If the teachers' soft labels equal the gold
    labels on the training split, the lever is text the teachers have not fitted. Three runs, each with
    its predictions written down beforehand, then measured the gain, its growth with the amount of text,
    and its survival under the matched-compute and second-student controls.
\end{enumerate}
The title changed with the evidence. The registered title named a dynamic, multi-teacher, homogeneous
and robust framework. The grid measured each of those four properties and found none of them to be the
source of the gain, so the final title names what is.
